\begingroup
\renewcommand{\arraystretch}{1.45}
\Needspace{32\baselineskip}
\InterpolationArticle{35}
Wenn die Funktion $X$ nicht mit $\sin nx$ abbricht, sondern weiterläuft und die folgenden Terme durch $\delta^{\mu+1}\sin(\mu+1)x+\delta^{\mu+2}\sin(\mu+2)x$ usw. ausgedrückt werden, so werden wir haben:
\[
\begin{aligned}
&\zeta'={}\frac{1}{\sin\mu a}\{\delta'\sin\mu a+\delta^{\mu+1}\sin 2\mu a+\delta^{2\mu+1}\sin 3\mu a+\text{usw.}\\[5pt]
&\qquad+\delta^{2\mu-1}\sin\mu a+\delta^{3\mu-1}\sin 2\mu a+\delta^{4\mu-1}\sin 3\mu a+\text{usw.}\}\\[5pt]
&\zeta''={}\frac{1}{\sin\mu a}\{\delta''\sin\mu a+\delta^{\mu+2}\sin 2\mu a+\delta^{2\mu+2}\sin 3\mu a+\text{usw.}\\[5pt]
&\qquad+\delta^{2\mu-2}\sin\mu a+\delta^{3\mu-2}\sin 2\mu a+\delta^{4\mu-2}\sin 3\mu a+\text{usw.}\}\\[5pt]
&\zeta'''={}\frac{1}{\sin\mu a}\{\delta'''\sin\mu a+\delta^{\mu+3}\sin 2\mu a+\delta^{2\mu+3}\sin 3\mu a+\text{usw.}\\[5pt]
&\qquad+\delta^{2\mu-3}\sin\mu a+\delta^{3\mu-3}\sin 2\mu a+\delta^{4\mu-3}\sin 3\mu a+\text{usw.}\}
\end{aligned}
\]
usw. Für den letzten Koeffizienten dagegen
\[
\begin{aligned}
&\zeta^{\mu}=\frac{1}{\sin\mu a}(\delta^{\mu}\sin\mu a+\delta^{2\mu}\sin 2\mu a+\delta^{3\mu}\sin 3\mu a+\text{usw.})
\end{aligned}
\]
\Needspace{40\baselineskip}
Diese Formeln zeigen, inwieweit der Unterschied zwischen $X$ und $X''$ vernachlässigt werden kann. Diese letztere Formel wird unter allen ähnlichen, die den $\mu$ vorgelegten Werten genügen, die einfachste sein; alle diese werden aber in der Formel $X''+Y\sin x(\cos\mu x-\cos\mu a)$ enthalten sein, wobei $Y$ wie in Artikel 32 eine unbestimmte Funktion des Bogens $x$ bezeichnet, die frei von Sinusgliedern ist. Und allgemein: Ist $X'''$ irgendeine Funktion derselben Form wie $X$, d. h. nur Sinusglieder enthaltend, die den $\mu$ gegebenen Werten genügt, so wird die Formel $X'''+Y\sin x(\cos\mu x-\cos\mu a)$ alle derartigen Funktionen enthalten; bei geeigneter Bestimmung von $Y$ kann sie auf die $\mu^{\mathrm{te}}$ Ordnung herabgesetzt werden, wobei notwendig die Funktion $X''$ selbst hervorgehen muss. Ganz ähnlich wie in Artikel 32 ergibt sich hierzu folgende allgemeine Regel: Für jeden Term in $X'''$ von der Form $L\sin\lambda x$, wobei $\lambda$ größer als $\mu$ ist, muss man in $X''$, indem man $\lambda=k\mu+\lambda'$ setzt, sodass $k\mu$ dasjenige Vielfache von $\mu$ ist, das unmittelbar unter $\lambda$ liegt, und somit $\lambda'$ zwischen den Grenzen $1$ und $\mu$ einschließlich liegt, die Terme
\[
\begin{aligned}
&L\frac{\sin(k+1)\mu a}{\sin\mu a}\sin\lambda'x+L\frac{\sin k\mu a}{\sin\mu a}\sin(\mu-\lambda')x
\end{aligned}
\]
einsetzen, die sich, sooft $\lambda'=\mu$ wird, auf den einzigen Term $L\frac{\sin\lambda a}{\sin\mu a}\sin\mu x$ reduzieren.
\SourceLeaf{327}
\Needspace{25\baselineskip}
\InterpolationArticle{36}
Die im vorigen Artikel und in Artikel 32 angegebenen Transformationen lassen sich übersichtlicher aus einem gewissen allgemeinen Satz ableiten, den wir hier noch kurz berühren wollen, da er auch für sich recht elegant ist.

\textsc{Satz.} \textit{Es bezeichnen $\lambda,\lambda',\lambda''$ beliebige ganze Zahlen. Ferner bezeichne $\mu$ eine ganze Zahl, welche die Unterschiede zwischen ihnen, $\lambda'-\lambda,\lambda''-\lambda',\lambda-\lambda''$, teilt (z. B. die Einheit), $x$ einen unbestimmten Bogen und $a$ einen bestimmten Bogen: Die Funktionen}
\[
\begin{aligned}
&P=\sin(\lambda'-\lambda'')a\cos\lambda x+\sin(\lambda''-\lambda)a\cos\lambda'x+\sin(\lambda-\lambda')a\cos\lambda''x\\[5pt]
&Q=\sin(\lambda'-\lambda'')a\sin\lambda x+\sin(\lambda''-\lambda)a\sin\lambda'x+\sin(\lambda-\lambda')a\cos\lambda''x
\end{aligned}
\]
\textit{werden durch} $\cos\mu x-\cos\mu a$ \textit{teilbar sein.}

\Needspace{29\baselineskip}
\textit{Beweis.} Wenn $\lambda$ durch $\mu$ teilbar ist, sodass auch $\lambda',\lambda''$ durch $\mu$ teilbar sein werden, bestätigt sich leicht, dass der Wert von $P,Q$ bei der Einsetzung $x=a$ identisch $=0$ ist; daher wird $P$ durch die folgenden Ersetzungen nicht geändert:
\[
\begin{array}{lll}
\cos\lambda x&\text{wird ersetzt durch}&\cos\lambda x-\cos\lambda a\\
\cos\lambda'x&&\cos\lambda'x-\cos\lambda'a\\
\cos\lambda''x&&\cos\lambda''x-\cos\lambda''a
\end{array}
\]
und ebenso wenig $Q$ durch die Ersetzungen:
\[
\begin{array}{lll}
\sin\lambda x&\text{wird ersetzt durch}&\displaystyle\sin\lambda x-\frac{\sin\lambda a\sin\mu x}{\sin\mu a}\\[6pt]
\sin\lambda'x&&\displaystyle\sin\lambda'x-\frac{\sin\lambda'a\sin\mu x}{\sin\mu a}\\[6pt]
\sin\lambda''x&&\displaystyle\sin\lambda''x-\frac{\sin\lambda''a\sin\mu x}{\sin\mu a}
\end{array}
\]
\Needspace{32\baselineskip}
Diese sechs Ausdrücke sind aber durch $\cos\mu x-\cos\mu a$ teilbar; dies für den ersten und vierten bei einem positiven Wert von $\lambda$ gezeigt zu haben, genügt. Durch Multiplikation bestätigt sich nämlich leicht, dass
\[
\begin{aligned}
&\sin\mu a(\cos\lambda x-\cos\lambda a)\\[5pt]
&=(\cos\mu x-\cos\mu a)\{2\sin\mu a\cos(\lambda-\mu)x+2\sin 2\mu a\cos(\lambda-2\mu)x\\[5pt]
&\quad \qquad+2\sin 3\mu a\cos(\lambda-3\mu)x+\text{usw.}+2\sin(\lambda-\mu)a\cos\mu x+\sin\lambda a\}
\end{aligned}
\]
\[
\begin{aligned}
&\sin\mu a\sin\lambda x-\sin\lambda a\sin\mu x\\[5pt]
&=(\cos\mu x-\cos\mu a)\{2\sin\mu a\sin(\lambda-\mu)x+2\sin 2\mu a\sin(\lambda-2\mu)x\\[5pt]
&\quad \qquad+2\sin 3\mu a\sin(\lambda-3\mu)x+\text{usw.}+2\sin(\lambda-\mu)a\sin\mu x\}
\end{aligned}
\]
Der Fall, in dem $\lambda$ negativ ist, führt von selbst auf diesen zurück. Hieraus ist ersichtlich, dass die Funktionen
\SourceLeaf{328}
$P$, $Q$ aus durch $\cos\mu x-\cos\mu a$ teilbaren Teilen zusammengesetzt und daher auch selbst durch diesen Divisor teilbar sind.

\Needspace{34\baselineskip}
II. Wenn $\lambda$ nicht durch $\mu$ teilbar ist, sei $l$ eine beliebige ganze, durch $\mu$ teilbare Zahl, und man setze $\lambda=l+\theta$, $\lambda'=l'+\theta$, $\lambda''=l''+\theta$; dann werden auch $l'$, $l''$ durch $\mu$ teilbar sein. Nun ist ersichtlich, dass, wenn man
\[
\begin{aligned}
&\sin(l'-l'')a\cos lx+\sin(l''-l)a\cos l'x+\sin(l-l')a\cos l''x=P'\\[5pt]
&\sin(l'-l'')a\sin lx+\sin(l''-l)a\sin l'x+\sin(l-l')a\sin l''x=Q'
\end{aligned}
\]
setzt,
\[
\begin{aligned}
&P=P'\cos\theta x-\theta'\sin\theta x,\quad Q=P'\sin\theta x+Q'\cos\theta x,
\end{aligned}
\]
wird und dass die Funktionen $P'$, $Q'$, da sie zum bereits erledigten ersten Fall gehören, durch $\cos\mu x-\cos\mu a$ teilbar sind: Daher werden offensichtlich auch $P$ und $Q$ durch $\cos\mu x-\cos\mu a$ teilbar sein. Q.~E.~D. Im Übrigen ist der Beweis des ersten Falles so geführt, dass er, sooft $\sin\mu a=0$ ist, nicht ohne gewisse Erläuterungen angewendet werden kann; dann wird aber $P=0$, $Q=0$, sodass es überhaupt keines Beweises bedarf.

\Needspace{24\baselineskip}
Setzt man daher $\lambda-\lambda'=k\mu$, $\lambda-\lambda''=(k+1)\mu$, so sind ersichtlich durch $\cos\mu x-\cos\mu a$ teilbar
\[
\begin{aligned}
&\sin\mu a\cos\lambda x-\sin(k+1)\mu a\cos\lambda'x+\sin k\mu a\cos(\lambda'-\mu)x\\[5pt]
&\sin\mu a\sin\lambda x-\sin(k+1)\mu a\sin\lambda'x+\sin k\mu a\sin(\lambda'-\mu)x
\end{aligned}
\]
und somit auch
\[
\begin{aligned}
&\displaystyle\cos\lambda x-\frac{\sin(k+1)\mu a\cos\lambda'x-\sin k\mu a\cos(\mu-\lambda')x}{\sin\mu a}\\[5pt]
&\displaystyle\sin\lambda x-\frac{\sin(k+1)\mu a\sin\lambda'x+\sin k\mu a\sin(\mu-\lambda')x}{\sin\mu a}
\end{aligned}
\]
woraus der Grund für die Substitutionen in den Artikeln 32, 35 sogleich hervorleuchtet: Denn der Quotient der letzteren Division wird nur aus Sinusgliedern bestehen und somit offensichtlich wiederum durch $\sin x$ teilbar sein.

\endgroup
